\begin{figure}[t]
    \centering 
    \label{fig:cobertura_fluxos}
    \caption{Cobertura de fluxos por Caso de Uso na Matriz de Rastreabilidade}
    \vspace{-1em}
    \resizebox{\textwidth}{!}{
    \begin{tikzpicture}
        \begin{axis}[
            xbar stacked, 
            width=14cm, 
            height=7cm, 
            enlarge y limits=0.15, % Dá um respiro maior nas margens superior e inferior
            xmin=0, 
            xlabel={Quantidade de Fluxos},
            ylabel={Casos de Uso},
            ylabel style={xshift=-10pt}, % Empurra o texto do eixo Y mais para a esquerda
            symbolic y coords={CdU10, CdU03, CdU02, CdU01},
            ytick=data, 
            % --- MELHORIAS ESTÉTICAS ---
            axis x line*=bottom, % Deixa apenas a linha do eixo X embaixo
            axis y line*=left,   % Deixa apenas a linha do eixo Y à esquerda
            xmajorgrids=true,    % Ativa a grade no eixo X
            grid style={dashed, gray!30}, % Deixa a grade tracejada e discreta
            bar width=18pt,      % Afina as barras para ficarem mais elegantes
            % Legenda sem borda e com mais espaço entre as palavras
            legend style={at={(0.5,-0.19)}, anchor=north, legend columns=-1, draw=none, /tikz/every even column/.append style={column sep=0.5cm}}, 
            % Código que imprime o número na barra apenas se ele for maior que zero
            nodes near coords={
                \pgfmathfloatifflags{\pgfplotspointmeta}{0}{}{\pgfmathprintnumber{\pgfplotspointmeta}}
            },
            nodes near coords align={center}, 
            % Formatação do número dentro da barra (negrito e tamanho menor)
            every node near coord/.append style={font=\footnotesize\bfseries, text=black!80},
        ]
        
        % 1ª SÉRIE: Fluxos 100% Cobertos (Verde com borda mais escura)
        \addplot[fill=mutedDGreen, draw=mutedDGreen!80!black, thick] coordinates {
            (2,CdU01) 
            (3,CdU02) 
            (1,CdU03)
            (1,CdU10)
        };
        
        % 2ª SÉRIE: Fluxos Parcialmente Cobertos (Laranja com borda mais escura)
        \addplot[fill=mutedOrange, draw=mutedOrange!80!black, thick] coordinates {
            (0,CdU01) % Este zero não será impresso, deixando o gráfico limpo!
            (3,CdU02) 
            (2,CdU03)
            (2,CdU10)
        };

        % 3ª SÉRIE: Fluxos Não Cobertos (Vermelho suave com borda)
        \addplot[fill=mutedRed, draw=mutedRed!80!black, thick] coordinates { 
            (2,CdU01) 
            (1,CdU02) 
            (1,CdU03)
            (4,CdU10)
        };
        
        \legend{Totalmente Cobertos, Parcialmente Cobertos, Não Cobertos}
        
        \end{axis}
    \end{tikzpicture}
    }
    {\footnotesize \textbf{Fonte:} Elaborado pelo autor.}
\end{figure}